\documentclass[10pt,a4paper,twoside]{article}

\usepackage[utf8]{inputenc}
\usepackage[brazil,portuges]{babel}
\usepackage{csquotes}

\usepackage[
  backend=biber,
  style=abnt,
  language=auto,
]{biblatex}

  \addbibresource{../refs/referencias.bib}

\usepackage{soul}
\usepackage{xcolor}
\usepackage{framed}

\usepackage{kpfonts}
\usepackage[T1]{fontenc}

\usepackage{a4wide}

  \setlength{\headheight}{12.5pt}

\usepackage{tabularx}

  \newcolumntype{C}{>{\centering\arraybackslash}X}

\usepackage{graphicx}
\usepackage{multirow}
\usepackage{multicol}
\usepackage{hhline}
\usepackage{booktabs}
\usepackage[inline]{enumitem}
\usepackage{xfrac}
\usepackage{mathtools}

  \newcommand{\Binary}{\ensuremath{\mathbb{B}}}
  \newcommand{\Natural}{\ensuremath{\mathbb{N}}}
  \newcommand{\Integer}{\ensuremath{\mathbb{Z}}}
  \newcommand{\Rational}{\ensuremath{\mathbb{Q}}}
  \newcommand{\Real}{\ensuremath{\mathbb{R}}}

  \DeclarePairedDelimiter{\floor}{\lfloor}{\rfloor}
  \DeclarePairedDelimiter{\ceil}{\lceil}{\rceil}
  \DeclarePairedDelimiter{\abs}{\lvert}{\rvert}

\usepackage{amsthm}

  \newtheorem{proposition}{Proposição}[section]
  \theoremstyle{definition}
  \newtheorem{definition}{Definição}[section]

\usepackage{multicol}

  \newcommand{\finishcols}{\vfill\vphantom{.}}

\usepackage[portuguese,algoruled,vlined,noend]{algorithm2e}

  \SetKwInput{Algoritmo}{Algoritmo}
  \SetKw{Retorne}{retorne}
  \SetKwFor{Para}{para}{faça}{fimpara}
  \SetKw{Escreva}{escreva}
  \DontPrintSemicolon

\usepackage{hyperref}
\usepackage[brazil]{varioref}
\usepackage[brazilian]{cleveref}

  \def\disciplinacodigotext{QXD0115}
  \def\disciplinanometext{Estrutura de Dados Avançada}
  \def\semestretext{2022.1}
  \def\provatype{Notas de Aula}
  \def\provatitle{Árvores AVL são balanceadas}
  \def\provatext{\provatype{} \ifx\provatype\empty {} \else {--} \fi \provatitle{}}
  \def\provacodtext{}
  \def\professorabbrv{Prof.}
  \def\professornometext{Arthur Araruna}

\usepackage{fancyhdr}

  \pagestyle{fancy}
  \fancyhead{}
  \lhead{\disciplinacodigotext}
  \chead{\provatext}
  \rhead{\semestretext}
  \fancyfoot{}
  \fancyfoot[LE,RO]{\colorbox{black!20}{\thepage}}

% * NOTE: Redefinição feita para evitar confusão com exponenciações.
\renewcommand{\thefootnote}{(\roman{footnote})}

\begin{document}
  \noindent
  \begin{minipage}{\textwidth}
    \begin{tabularx}{\textwidth}{cXc}
      \multirow{4}{*}{\includegraphics[height=1.8cm]{brasaoufc.pdf}} & {\bf UNIVERSIDADE FEDERAL DO CEARÁ}         & \multirow{3}{*}{} \\
                                                                     & {\bf Campus de Quixadá}                     &                   \\
                                                                     & \professorabbrv\enspace\professornometext   &                   \\
                                                                     & \disciplinacodigotext - \disciplinanometext & {}                \\
    \end{tabularx}
  \end{minipage}

  \bigskip
  \bigskip

  \part*{\provatitle}

    \section{Conceitos}

     Vamos relembrar o nossos conceitos que serão necessários para o resultado a ser demonstrado:

     \begin{definition}
       Uma árvore binária \( T \), com \( n \) nós, tal que \( h(T) = O(\lg n) \) é dita \emph{balanceada}.
     \end{definition}

     \begin{definition}
       O \emph{balanço} de um nó binário \( n \) é definido pela função \( bal(n) \coloneqq h(n.dir) - h(n.esq) \).
     \end{definition}

     \begin{definition}
       Uma árvore binária \( T \) é \emph{AVL} quando, para todo nó \( n \in T \), \( \abs{bal(n)} \leq 1 \).
     \end{definition}

    \section{Demonstração}

     \begin{framed}
       {\bfseries Observação:}

       A argumentação abaixo não é baseada na que consta em~\textcite{Szwarcfiter2010} pois uma das afirmações dadas como válidas não é de fácil demonstração.
     \end{framed}

     Note que, se mostrarmos que a altura de uma árvore AVL \ul{qualquer} sempre é dada no máximo por uma função logarítmica do seu tamanho, poderemos concluir que a ordem de grandeza da altura não pode ser maior que logarítmica. Isso é exatamente o que significa, considerando \( h \) sendo a altura e \( n \) sendo o tamanho, que \( h = O(\lg n) \) e, com isso, que toda AVL seria balanceada.

     Nesse sentido, uma das formas possíveis é se conseguirmos construir uma desigualdade da forma \( h \leq c \cdot k + \lg n \) para constantes \( c \) e \( k \), com \( c > 0 \). A estratégia que seguiremos será um pouco mais restrita, porém permitirá concluirmos a mesma coisa.

     Em vez de analisarmos uma árvore qualquer, vamos observar, para um dado \( h \) arbitrário, uma árvore AVL com altura \( h \) que contenha o mínimo possível de nós. Se essa árvore respeitar a desigualdade anterior, poderemos concluir que \emph{todas} as AVL com essa mesma altura também a respeitarão\footnote{Pois já que para essa tal árvore ``mínima'' o lado direito da desigualdade é o menor possível, todas as outras apresentarão valores ainda maiores, o que mantém a desigualdade.}.

    \subsection{Estrutura da árvore}

      Para entendermos como construir uma árvore como desejamos, observe o raciocínio a seguir.

      Imagine que \( T \) seja uma árvore AVL com altura \( h \), e considere a situação em que encontramos um nó \( n \in T \) com \( bal(n) = 0 \). Veja que se substituirmos, por exemplo, a sub-árvore \( n.esq \) por uma outra árvore AVL arbitrária com altura \( h(n.esq) - 1 \), estaremos construindo uma nova árvore \( T^\prime \) menor que \( T \) mas de mesma altura e ainda AVL. Veja a~\Vref{fig:troca-sarv}.

      \begin{figure}[hbt]
        \caption{\label{fig:troca-sarv}}
        \centering
        \begin{tikzpicture}[transform shape, scale=.9, nodes={draw, circle, minimum width=3.5ex}, sibling distance=5.5em, level distance=6ex]
        \end{tikzpicture}
      \end{figure}

      Isso nos permite concluir que a AVL ``mínima'' de altura \( h \) deve possuir todos os balanços não-nulos, caso contrário ainda poderíamos diminuir seu tamanho usando o raciocínio descrito.

    \subsection{Contagem de nós da árvore}
      \label{str:contagem-nos}

      Vamos definir a função \( N(h) \) de forma a representar o menor número de nós possível para uma árvore AVL de altura \( h \), que chamaremos de \( T \).

      Como vimos anteriormente, as sub-árvores de \( T \) devem possuir alturas \( h-1 \) e \( h-2 \), não necessariamente nesta ordem. Além disso, podemos afirmar que ambas se tratam de árvores AVL com o mínimo de nós para a respectiva altura, caso contrário não faria sentido \( T \) possuir tamanho mínimo (já que elas poderiam ser trocadas por outras árvores de mesma altura e menos nós).

      Dessa forma, podemos afirmar que o valor de \( N(h) \)\footnote{O caso \( h = 0 \) não é considerado, mas é trivial, dado que a única árvore com altura 0 é a árvore vazia, com 0 nós.} pode ser calculado como \[ N(h) = \begin{cases}
          N(h-1) + N(h-2) + 1, & \text{se } h \geq 2; \\
          2                  , & \text{se } h = 2;    \\
          1                  , & \text{se } h = 1.
        \end{cases} \]

      Observe que a função \( N(h) \) é monotonicamente crescente\footnote{Isso porque todos os valores envolvidos são não-negativos e as operações realizadas são apenas adições.}. Com isso, e observando os casos-base, podemos afirmar que \( N(h) \geq 2N(h-2) \)\footnote{Considere que substituímos o termo \( N(h-1) \) por um termo \( N(h-2) \) e o termo 1 por um termo 0. Em ambas as substituições, os novos valores são menores.} para todo \( h \geq 3 \).

    \section{Conclusão}

     Resolvendo\footnote{Veja argumentação no Apêndice.} a nova equação de recorrência obtida, obtemos que \( N(h) \geq 2^{\frac{h-1}{2}} \). Lembrando que \( N(h) \) é o número de nós de \( T \), que passaremos a denominar por \( n \), concluímos que, para todo \( h \geq 3 \),
     \begin{alignat*}{3}
       n = N(h) & \geq 2^{\frac{h-1}{2}} & \quad &  & \implies \\
       \lg n    & \geq \frac{h-1}{2}     & \quad &  & \implies \\
       h        & \leq 1 + 2 \lg n       & \quad &  &
     \end{alignat*}

     Dessa forma, temos todas as informações necessárias para concluirmos que \( h = O(\lg n) \).

  \part*{Apêndice}

    \section{Resolução da Recorrência}

     Na~\Vref{str:contagem-nos} determinamos que \( N(h) \geq 2N(h-2) \) para todo \( h \geq 3 \). Vamos agora mostrar que isso implica em \( N(h) \geq 2^{\frac{h-1}{2}} \) para todo \( h \geq 3 \).

     Observe a seguinte sequência de desigualdades, obtida por sucessivas aplicações da desigualdade mencionada acima.

     \begin{alignat*}{2}
       N(h) & \geq 2N(h-2)           &  & = 2^1 N(h-2)  \\
            & \geq 2(2N(h-4))        &  & = 2^2 N(h-4)  \\
            & \geq 2^2(2N(h-6))      &  & = 2^3 N(h-6)  \\
            & \geq 2^3(2N(h-8))      &  & = 2^4 N(h-8)  \\
            & \dots                  &  & \dots         \\
            & \geq 2^{k-1}(2N(h-2k)) &  & = 2^k N(h-2k)
     \end{alignat*}

     Considerando que a última substituição nos tenha trazido a um caso base, temos as seguintes possibilidades:

     \begin{description}
       \item[CASO \( h-2k = 2 \):] podemos concluir que \( k = \frac{h-2}{2} \) e que \[ N(h) \geq 2^{\frac{h-2}{2}} N(2) = 2^{\frac{h-2}{2}} \cdot 2 = 2^{\frac{h}{2}} \geq 2^{\frac{h-1}{2}}. \]
       \item[CASO \( h-2k = 1 \):] podemos concluir que \( k = \frac{h-1}{2} \) e que \[ N(h) \geq 2^{\frac{h-1}{2}} N(1) = 2^{\frac{h-1}{2}} \cdot 1 = 2^{\frac{h-1}{2}}. \]
     \end{description}

     \printbibliography

\end{document}